\documentclass{article}

\usepackage{float}
\usepackage[backref=page]{hyperref}

%----------Title----------
\title{\textbf{Diary-2022}}
\author{Fernando Gutiérrez-Canales}
\date{}

\begin{document}
\maketitle

\section{11 de juin de 2022}

Today I went to the \( 39^{e}  \) \textit{marche de poésie} at \textit{la place saint-sulpice paris \( 6 ^{e} \)}. The whole place was full of people and books. I like very much that this kind of book fairs take place outdoors. The impressive Saint-Sulpice church seemed to be watching the whole crew the whole time. I steped inside and the first thing I saw was the most beautiful girl in France, followed by the most beautiful girl in France and then I saw the most beautiful girl in France and so on... All the girls in the place were so beautiful that we, the men, seemed to be even uglier than we really are. I saw all these girls going around and following poets and writers that looked suspiciouly like Michel Houllebecq. Not a single Frederic Beigbeder in town, just Houllebecqs all around.\\

Everyone was smoking and drinking wine and eating different types of cheese. There was a little stage where authors read aloud its poetry. When I looked at the stage to see what was happening, I saw three men and one woman sitting in front of the public, small-talking between each other. They were the authors. Out of the three authors, one of the men seemed really young, while the others two looked like they were in their 50s already. One of the older poets had a spanish name and a nice chubby face. The other one seemed like an american elementary school teacher taking a free weekend in Paris. The young man seemed Danish or Swedish or German or Austrian. He looked like a hacker, but I couldn't explain exactly why. \\

Out of the three authors, the one that impressed me the most was the woman. She looked like a really depressed person. She was wearing a pair of jeans with sneakers and a pretty big and used blue T-shirt. Her hair was a mess and she wasn't talking too much with the other authors. However, when it was her turn to read her poems in front of everybody, she took a deeph breath and then started to read a beautiful poem called \textit{la tristesse} (The sadness). In her poem, she talked about how the sadness was unbearably heavy, how she cannot longer hold it anymore, how her mind is inside a fog. She talked about Leonard Cohen, about the fear for the future, which is gigantic. She talked about her addictions, about her ugliness, her sexual desire, her lack of sexual desire, her voice, her right hand and her left foot. She talked about the immense hate she feels for herself. She talked to everyone of us about her most hidden fear. She talked about abortion, she talked about burned-to death kids. She talked about the time she saw a dying cat, with its guts out. She talked about the guilt. She talked about the unbearably unbearavly unbearably unbearably weight of guilt. She talked about hate. She talked about painful masturbation. She talked about kidney stones. She talked about how she does miss her aunt. She talked about how she hates society, how she loves society. She talked about the monster. She talked about perverts. She talked about me. She is me. She talked about all the monstrous lies she had been told her entire life to everyone else around her. She talked about being defeated by everyone and everything. She talked about being knocked unconscious by stress and fear and anger and axiety and the expectation of others. She talked about what expectations are placed on her, about what she thinks other people expect of her and how that is totally destroying her. She talked about suicide. She talked about being loved. She talked about being kind to other people. She talked about how that is impossible for her, a woman who is a snake. She talked about the indifference she feels towards anyone but herself. She talked about \textit{La malagueña}. She talked about \textit{musique de chambre}, about Glenn Gould, about Schutz. She talked about Chopin, about Martha Arguerich playing Chopin. She talked about true happiness, about how true happiness is possible in this world because Martha Arguerich plays Chopin and Glenn Gould plays the Goldberg variations. She talked about Schutz' \textit{Die sieben Worte Jesu Christi am Kreuz}, about how is the most beautiful thing in the world. She talked about God, the Bhagavad Gita, the Torá and Mohammed. She talked about Horace and Kavafis. She talked about Georg Trakl, a heart wrapped in spiders. She talked about Paul Celan's suicide, black milk and snakes in the morning. She talked about Jorge Cuesta and Origen Adamantius. She talked about Philipp Mainländer. She talked about Mayakowski. She talked about Kenji Miyazawa. She talked about Idea Vilariño, the most beautiful poems ever written. She talked about Pizarnik and Plath, poor little girls. She talked about, she talked about, she talked about. She talked about Anne Sexton, who is here with me. She talked about the lonely and sad Pavese. She talked about Ida Vitale. She talked about Inés Arredondo. She talked about Lucian Blaga. She talked about poetry, how poetry fills the world, how poetry is the only thing in this world. She talked to us. She talked to us. She talked to us. She talked to us. \\

But I am and always be someone stupid, who never pays attention when recquired and never get anything done by himself. I didn't understand a single word.

\end{document}

